\section{Discussion}
\label{sec:discussion}

\subsection{What the result means for practice}

The finding that sensitivity tracks the task more than the judge has a direct
consequence: choosing a stronger judge is not a reliable way to buy stability.
Across the three judges measured, the most and least sensitive judge on coherence
differ by $0.044$ $\Delta\mathrm{JSS}$, while the most and least sensitive
\emph{task} within a single judge differ by $0.162$. A practitioner worried about
reproducibility gains more by reconsidering what they are asking the judge to do
than by upgrading the model.

The ordinal task is the least stable of the four, and we think the mechanism is
general rather than particular to coherence. A five-point scale gives a rewording
more ways to move the answer, and the weighted kappa shows most of that movement
is a single scale point. Evaluation designs that collect a graded score and then
threshold it are absorbing this instability silently; a binary or pairwise
formulation of the same underlying question is measurably more reproducible.

\subsection{Refusal is an evaluation failure mode, not an aside}

One judge declined a quarter of the relevance items while another, on identical
prompts, declined none. For anyone using an LLM judge over medical or otherwise
sensitive corpora this is a first-order problem: the missing verdicts are not
missing at random, and an evaluation harness that silently treats a decline as a
parse failure will report a format-adherence problem that does not exist.

It also breaks a convenience that most analyses assume. Because a rewording can
itself change whether a judge answers, refusal is an outcome of the manipulation
rather than a nuisance to be dropped, and conditioning on the answered subset is
selection on a post-treatment variable. We report worst-case bounds alongside the
conditioned estimate for this reason. Where refusals are common those bounds are
wide, and we would rather show a wide identified region than a precise number
that depends on an assumption nobody can check.

\subsection{On building a benchmark that can be wrong}

The defects we found in our own rebuild are, we think, the most transferable part
of this work. Each one passed every check that existed at the time.

A deterministic rotation assigning template pairs to items ran in lockstep with
the alternating emission of correct and incorrect answers, making the template
pair a perfect answer key --- and skewing two templates' label balance to
$75/25$, so a constant-answer judge would have appeared to have a large
template preference. Candidate responses reused across pairings carried
contradictory gold labels, so a judge reasoning correctly and consistently was
scored wrong on nearly half the preference items. A length-balancing step that
correctly removed a verbosity shortcut cut only one of the two buckets, and cut
it by vote margin, leaving the buckets matched on length and far apart on
annotator agreement.

None of these is detectable by inspecting labels, provenance, or aggregate
statistics. All three were found by asking what a heuristic with no understanding
would score, and by checking whether the artifact contained a rule that predicted
the answer. We would encourage that as routine practice: before reporting what
models score on a benchmark, report what nothing scores on it, and report it
with an interval and a support count rather than a bare number.

\subsection{Limitations}

The evidence base is three judges from one vendor family. The task ordering is
consistent across them, but that consistency is a within-family observation and
we do not claim it generalises; the harness supports fourteen judges across seven
providers and eleven are unrun. Two of the twelve judge--task cells report no
endpoint, their support having fallen below the pre-registered floor. One judge
rejects an explicit temperature parameter and therefore ran at a provider
default while the others ran at zero; this is recorded per call rather than
inferred, and that judge's cells should be read accordingly.

The design measures sensitivity within an item across a template pair, and cannot
compare templates to one another: no item appears under more than one pair, so
template identity is nested inside item identity and the two effects are not
separable. Our equivalence checks bound specific failure modes --- label space,
polarity, construct, non-triviality --- but nothing computable from the template
strings establishes that a competent reader maps two wordings onto the same
question. Semantic drift is reduced as an alternative explanation, not
eliminated.

Finally, the effect sizes are modest in absolute terms. A $\Delta\mathrm{JSS}$ of
$-0.04$ on factuality means a judge that reproduces itself on the identical
prompt loses four points of agreement under rewording. Whether that matters
depends on the decision being made downstream, and we have declared in advance
that we do not regard anything below $0.02$ as practically meaningful.

\section{Conclusion}

We set out to measure whether LLM judges return the same verdict when asked the
same question in different words, and to do so on a benchmark that could not be
passed without judging. Across three judges and four tasks, every measurable cell
shows a real loss of agreement under meaning-preserving rewording, beyond what
the judge's own decoding variance explains. The size of that loss depends more on
the evaluation task than on the model performing it.

The benchmark, the analysis plan fixed before the sweep, the raw judge outputs,
and the heuristic baselines that would beat the tasks are all released. So are
the defects we found in our own construction, and the measurements that exposed
them.
